% Source-derived chapter 03: Geometric constructions
% Original source: motivic-invariants-birational-maps
\chapter{Geometric constructions}
\label{motivic-invariants-birational-maps-chapter-03}
\source{motivic-invariants-birational-maps}

\begin{remark}[Source section]
\label{motivic-invariants-birational-maps-chapter-03-source}
\source{motivic-invariants-birational-maps}
\source{motivic-invariants-birational-maps}
This chapter reproduces the corresponding section of the retrieved
LaTeX source, including its definitions, statements, examples, and
proofs. It has not yet been translated into Lean.
\notready
\end{remark}

% The source section title was: Geometric constructions
\label{sec:geometric-constructions}
\source{motivic-invariants-birational-maps}



\subsection{Strong birational isomorphisms}
\hfill

Recall that two varieties $X$ and $Y$ are isomorphic
in codimension one if there exists a birational
map $\gamma: X \dto Y$ such that
$\Ex(\gamma)$ and $\Ex(\gamma^{-1})$ have codimension $\ge 2$.
We introduce a motivic version of this notion:

\begin{definition}
\source{motivic-invariants-birational-maps}
\notready\label{def:strong-birat}
\source{motivic-invariants-birational-maps}
We say that two varieties
$X$ and $Y$ are {\it strongly
birational}
if there exists $\gamma \in \Bir(X,Y)$
such that 
$c(\gamma) = 0$.
We say that $X$ is 
{\it strongly rational}
if $X$ is strongly birational to $\P^n$.
\end{definition}

Note that strong birationality
is an equivalence relation by the
additivity of $c$ (Lemma \ref{lem-homc}).
Furthermore, isomorphism in codimension one
implies strong birationality.


\medskip

When $\k$ is a field of characteristic zero,
it follows from Proposition
\ref{prop:Picnb} that strongly birational varieties
have equal ranks of their
geometric N\'eron--Severi groups.
In particular, a strongly rational
variety has geometric N\'eron--Severi group 
 of
rank one. 

The following lemma provides some examples 
of strongly rational varieties.

\begin{lemma}
\source{motivic-invariants-birational-maps}
\notready\label{lem:GP-strongly-rational}
\source{motivic-invariants-birational-maps}
If $X$ is a smooth projective variety
containing an irreducible rational
divisor $D \subset X$
such that $U := X \setminus D$
is isomorphic to $\A^n$,
then $X$ is strongly rational.
In particular, if $G$ is a split reductive
algebraic group over a field $\k$,
then every 
homogeneous space
$X = G/P$
of Picard rank $\rho(X) = 1$
is strongly rational.
\end{lemma}
\begin{proof}
Consider the birational map 
$\phi :  X \dto \P^n$
induced by the isomorphism $U \simeq \A^n$.
We have
\[
\phi: X \hlto U \simeq \A^n \hto \P^n,
\]
so $c(\phi) = [\P^{n-1}] - [D] = 0$.
Thus $X$ is strongly rational.

Now we prove the second statement. 
A homogeneous space $X$ admits
the Bruhat stratification
into affine spaces \cite[Proposition 1.3]{Kock}. 
It has a dense open stratum isomorphic to $\A^n$ and 
the number of codimension one strata 
$\A^{n-1}$ in the Bruhat stratification
is equal to the Picard rank $\rho(X) = 1$ by~\cite[Example 1.9.1]{Fulton}.
The closure of the
$\A^{n-1}$ stratum  
is therefore the complement of the dense stratum $\A^n$.
It follows from the first part of the lemma that $X$
is strongly rational.
\end{proof}





\subsection{L-links}
\hfill


Let $X$, $Y$ be 
irreducible geometrically reduced
varieties.

\begin{Def} 
An \emph{L-link} with centers $X$ and $Y$
is a diagram 
\begin{equation}\label{eq:L-link}
\source{motivic-invariants-birational-maps}
 \xymatrix{
    & & T \ar[dl]_{\Bl_{X}} \ar[dr]^{\Bl_{Y}} & & \\
    X \ar@{^{(}->}[r]^{\text{closed}} & \XX \ar@{-->}[rr]^\psi &  & \XX' & \ar@{_{(}->}[l]_{\text{closed}} Y\\
}
\end{equation}
of blow ups with centers $X$ and $Y$,
subject to the following conditions:
\begin{enumerate}
    \item[(L1)] $\XX$ and $\XX'$ are 
 varieties smooth in codimension one,
 such that $[\XX] = [\XX']$ in $\Kt_0(\Var/\k)$.
    \item[(L2)] 
    $\XX$ and $\XX'$ are strongly birational.
    \item[(L3)]
    The exceptional divisors of
    $\Bl_{X}$ and $\Bl_{Y}$ are irreducible.
\end{enumerate}
\end{Def}

We say that the L-link $\psi$~\eqref{eq:L-link} is smooth 
(resp. projective)
if $X$, $Y$, $\cX$, $\cX'$, $T$
are smooth (resp. projective).
An L-link $\psi$
is called \emph{motivically non-trivial} if $c(\psi) \ne 0$.
For instance, this is the case when 
the link is smooth and
$X$ is not stably birational to $Y$.


\begin{example}
\source{motivic-invariants-birational-maps}
\notready
Special Cremona transformations~\cite{Crauder-Katz} give rise
to examples of L-links with
$\cX = \cX' = \P^n$
and $\gamma = \Id$. 
\end{example}

L-links are closely related to L-equivalence.
Recall that smooth
projective connected
varieties
$X$, $Y$
are
called
L-equivalent
if $\L^d([X]-[Y]) = 0$
in $\Kt_0(\Var/\k)$
for some $d \ge 0$.

\begin{lemma}
\source{motivic-invariants-birational-maps}
\notready\label{lemma:L-link}
\source{motivic-invariants-birational-maps}
Given an L-link \eqref{eq:L-link},
let $m = \codim_X(\XX)$,
$m' = \codim_Y(\XX')$. 
\begin{enumerate}
\item If the link is smooth, we have 
$\L([\P^{m-2}][X] - [\P^{m'-2}][Y]) = 0$ in $\Kt_0(\Var/\k)$.
\item Let $\phi = \gamma^{-1} \psi 
\in \Bir(\XX)$, where $\gamma$ comes from Definition \ref{def:strong-birat}.
Then 
$$c(\phi) = c(\psi) = [\P^{m-1} \times X] - [\P^{m'-1} \times Y]$$
in $\Z[\Bir/\k]$.
In particular, if the link \eqref{eq:L-link} is motivically non-trivial
(e.g. if $X$ and $Y$ are not stably
birational), then $c(\phi) \ne 0$.
\end{enumerate}
\end{lemma}
\begin{proof}
(1) follows by computing the class
of $[T]$ in two ways using the two
blow up morphisms and (L1).
For (2), we have $c(\phi) = c(\psi)$ by
(L2). 
Since $\cX$ is smooth in codimension one
and $X$ is geometrically reduced,
there exists a closed subscheme $Z \subset \cX$ of codimension two
such that both $\cX \bss Z$ and $X \bss Z$ are smooth.
As the exceptional divisor $E$ of $\Bl_{X}$
is irreducible by (L3), 
it follows that $E$ is birational to $\P^{m-1} \times X$.
Similarly,
the exceptional divisor of $\Bl_{Y}$
is birational to $\P^{m'-1} \times Y$.
Thus by Lemma~\ref{lem-homc}, we have
$c(\phi) = c(\psi) =  [\P^{m-1} \times X] - [\P^{m'-1} \times Y]$.
\end{proof}


\subsection{Elliptic curves}\label{ssec-dim3}
\source{motivic-invariants-birational-maps}
\hfill

Let $C$ be a smooth projective
connected
curve and let $k \in \Z$. 
Let $\Jac^k(C)$ denote
the Jacobian of divisors of degree $k$ on $C$.
Each $\Jac^k(C)$ is a torsor under
the algebraic group
$\Jac^0(C)$. 
If $C$ is a curve of genus one,
then
$\Jac^k(C)$
are twisted forms of $C$
which are not necessarily isomorphic to it.
We concentrate on curves $C$
of genus one and degree five for which
the only interesting Jacobian is $\Jac^2(C)$ \cite{ShinderZhang}.

\begin{proposition}
\source{motivic-invariants-birational-maps}
\notready\label{prop:elliptic}
\source{motivic-invariants-birational-maps}
Let $C$ be a genus one curve 
having
a divisor of degree $5$,
and let $C' = \Jac^2(C)$.
We have the following smooth projective L-link:
\begin{equation}\label{eq:L-link-curves}
\source{motivic-invariants-birational-maps}
 \xymatrix{
    & & T \ar[dl]_{\Bl_{C}} \ar[dr]^{\Bl_{C'}} & & \\
    C \ar@{^{(}->}[r] & Q^3 \ar@{-->}[rr]^\psi &  & \P^3 & \ar@{_{(}->}[l] C'.\\
}
\end{equation}
Here $Q^3$ is a smooth quadric threefold.
\end{proposition}
\begin{proof}
By Riemann-Roch, the linear
system of a divisor
of degree $5$ on $C$ 
gives rise to a degree $5$ embedding $C \subset \P^4$,
and it is well-known that 
$C$ is the scheme-theoretic intersection 
of the quadrics containing it, see e.g. \cite[Proposition 4.2(ii)]{Fisher-quintics}. 
Thus the map given by the linear system $|2H - C|$
is resolved by blowing up $C$, and we obtain
the Crauder--Katz quadro-cubic transformation \cite[Theorem 2.2(ii)]{Crauder-Katz}
\begin{equation}\label{eq:L-link-curves-big}
\source{motivic-invariants-birational-maps}
 \xymatrix{
    & & T' \ar[dl]_{\Bl_{C}} \ar[dr]^{\Bl_{S}} & & \\
    C \ar@{^{(}->}[r] & \P^4 \ar@{-->}[rr]^{\psi'} &  & \P^4 & \ar@{_{(}->}[l] S \\
}
\end{equation}
where $S$ is the surface
parameterizing the secant lines of $C$. 
Thus $S$ is isomorphic to $\Sym^2(C)$ 
(cf Proof of \cite[Theorem 3.3(A)]{Crauder-Katz})
and
is embedded as a quintic ruled surface $\Sym^2(C) \to \Jac^2(C)$
into $\P^4$.

To obtain the L-link 
\eqref{eq:L-link-curves} we take a hyperplane
section of $\eqref{eq:L-link-curves-big}$ as follows.
Take any smooth quadric $Q^3 \subset \P^4$ containing $C$
and restrict $\psi'$ onto $Q^3$. 
The image $\psi'(Q^3)$
is a hyperplane $\P^3$, 
and we obtain a birational map $\psi: Q^3 \dto \P^3$ and a diagram
\eqref{eq:L-link-curves} with $C' = S \cap \P^3$.
This curve is a section of the projection $S \to \Jac^2(C)$,
hence $C' \simeq \Jac^2(C)$,
and $C' \subset \P^3$
is a degree five genus one curve.

We claim that $Q^3$ contains $\k$-lines. 
If $\k$ is a finite field, then this is true
for any smooth quadric in $\P^4$, as the quadratic form
in 5 variables over a finite field is isomorphic to a direct
sum of two copies of the hyperbolic plane and a one-dimensional
form, by the Chevalley--Warning theorem. 
Let us assume that $\k$ is infinite.
Under the map $\psi$, general lines on $Q^3$ correspond to general
conics in $\P^3$ which are five-secant to $C'$. Such conics
are parameterized by planes $\P^2 \subset \P^3$, and taking
a general 
plane $\P^2 \subset \P^3$
defined over $\k$ we get a $\k$-rational line on $Q^3$.


Finally \eqref{eq:L-link-curves}
satisfies the definition of L-link
because $[Q^3] = [\P^3]$ as $Q^3$ 
contains a line~\cite[Example 2.8]{KuznetsovShinder},
and
$Q^3$ 
is strongly rational
(see e.g. Lemma \ref{lem:GP-strongly-rational}).
\end{proof}


\begin{theorem}
\source{motivic-invariants-birational-maps}
\notready\label{thm:ell-curves}
\source{motivic-invariants-birational-maps}
Let $\k$ be a 
field
of one of the following types:
\begin{itemize}
    \item a number field,
    \item a function field over a field $\F$,
    where $\F$ can be
    \begin{itemize}
        \item an algebraically closed field, 
        \item a number field,
        \item a finite field. 
    \end{itemize}
\end{itemize}
Then we can choose a curve
$C$ such that \eqref{eq:L-link-curves} is motivically non-trivial.
Furthermore, 
the set $I$ of isomorphism classes of such curves has the cardinality of $\k$.
\end{theorem}

Here by a function field, we mean
a function field of a positive-dimensional
$\F$-variety.

\medskip

Before we prove the theorem, we need a preliminary result 
on the arithmetic of elliptic curves.

\begin{lemma}
\source{motivic-invariants-birational-maps}
\notready\cite{Lang-Tate, ClarkLacy}
\label{lem:torsors}
\source{motivic-invariants-birational-maps}
For a field $\k$ 
as in Theorem \ref{thm:ell-curves},
there exist infinitely many 
genus one
curves $C$
without rational
points but having 
a divisor of degree $5$,
and such that $j(\Jac^0(C)) \ne 1728$.
The set $I$ of isomorphism classes of such curves has the cardinality of $\k$.
\end{lemma}
\begin{proof}
If $\k$ is an infinite 
finitely generated field over $\Q$ or over a finite field, 
then infinitely many
such curves exist by~\cite[Theorem 1.10]{ClarkLacy};
the condition $j(\Jac^0(C)) \ne 1728$ is not explicitly
checked in \cite{ClarkLacy}, however the elliptic
curves they use to produce torsors indeed satisfy this condition \cite[Section 4]{ClarkLacy}.

The remaining case is
a function field $\k$ over an algebraically
closed field $\F$. 
We apply~\cite[Theorem 7]{Lang-Tate}.
First of all since $\k$ is Hilbertian~\cite[Proposition 12.3.3, Theorem 13.4.2]{MR2445111},
it has infinitely many cyclic extensions
of degree $5$ which are
linearly disjoint in $\ol{\k}$~\cite[Corollary 16.2.7]{MR2445111}. 
In particular, these cyclic extensions are non-isomorphic.
Next we take any ordinary elliptic
curve $E/\F$
with $j(E) \ne 1728$ in $\F$.
Since $\F$ is algebraically closed,
$E(\F)$ has non-trivial $5$-torsion.
The base change $E_\k$ of $E$ to $\k$
has also non-trivial $5$-torsion
and $j(E_\k) \ne 1728$ in $\k$.
The group $E(\k) / 5E(\k)$
is finite by~\cite[Remark 1.3, Proposition 1.4]{ClarkLacy}. 
We thus conclude by~\cite[Theorem 7]{Lang-Tate}
that there exists infinitely many $E_\k$-torsors of index $5$.

Note that among the fields we consider, $\k$ is uncountable only if $\k$ is the function field over an algebraically closed field $\F$.
In this case $\F$ and $\k$ have the same cardinality.
As the set of elliptic curves $E/\F$
considered above has the same cardinality as $\F$,
the last statement follows.
\end{proof}

\begin{proof}[Proof of Theorem \ref{thm:ell-curves}]
Take the curve $C$ defined
by Lemma \ref{lem:torsors}
and set $C' = \Jac^2(C)$.
By \cite[Lemma 2.7]{ShinderZhang}
we have $C \not\simeq C'$
(the cited result assumes $\chr(\k) = 0$,
however the proof works without this assumption
since
there is no $\gs \in \Aut(E)$ of order $4$ when $j \ne 1728$
for any field
\cite[Proposition A.1.2]{SilvermanAEC}).
As $C$ and $C'$ are non isomorphic smooth curves, they are
not birational.
Hence the L-link~\eqref{eq:L-link-curves} is motivically non-trivial.
\end{proof}


\begin{remark}
\source{motivic-invariants-birational-maps}
\notready
The L-link \eqref{eq:L-link-curves} appears in the classification
of Fano threefolds of rank two by Mori--Mukai \cite[2) on page 117]{Mori-Mukai}.
It is a very interesting question whether
there exist other
nontrivial smooth projective
L-links between strongly rational varieties in dimension $3$. The ultimate
question in this direction
is to fully describe the image $c(\Bir(\P^3))$
using the Sarkisov link decomposition.
\end{remark}

\begin{remark}
\source{motivic-invariants-birational-maps}
\notready
Let $X \to B$ be an elliptic surface
with a multisection of degree five over $\k = \C$.
Let $Y = \Jac^2(X/B)$ be the relative Jacobian
of degree two divisors.
Then spreading out the construction
of Theorem \ref{thm:ell-curves},
one can show that
$[X] - [Y] \in \Image(c_{\P^3 \times B})$.
See~\cite{ShinderZhang}
for a detailed analysis when $X$, $Y$
are not birational in the case
of elliptic K3 surfaces.
\end{remark}


\subsection{K3 surfaces of degree $12$} 
\hfill

\begin{theorem}
\source{motivic-invariants-birational-maps}
\notready\label{thm:K3}
\source{motivic-invariants-birational-maps}
Let $\k \subset \C$.
There exist motivically 
non-trivial L-links 
\begin{equation}\label{eq:L-link-K3}
\source{motivic-invariants-birational-maps}
 \xymatrix{
    & & T \ar[dl]_{\Bl_{S_0}} \ar[dr]^{\Bl_{S'_0}} & & \\
    S_0 \ar@{^{(}->}[r] & \P^4 \ar@{-->}[rr]^\psi &  & \P^4 & \ar@{_{(}->}[l] S'_0\\
}
\end{equation}
where $S_0$ and $S'_0$ are projective
surfaces such that the minimal models $S$ and $S'$ of their normalizations
are non-isomorphic
K3 surfaces of geometric Picard
rank one and degree $12$, 
and $S'_\C$ is the unique Fourier--Mukai partner of $S_\C$.
Furthermore, 
the set $I$ of isomorphism classes of the K3 surfaces $S$ occurring in these L-links has the cardinality of $\k$.
\end{theorem}

Theorem \ref{thm:K3}, over $\C$,
is the main result of Hassett and Lai~\cite{HassettLai}.
We start by summarizing their construction,
and then we explain how to extend it to an arbitrary
subfield $\k \subset \C$ using Lefschetz pencils.

\begin{theorem}[{\cite[Theorem 2.1, Theorem 3.1]{HassettLai}}]
\source{motivic-invariants-birational-maps}
\notready\label{thm-HL}
\source{motivic-invariants-birational-maps}
Let $(S,H)$ be a general polarized complex K3 surface with $H^2 = 12$
and let $\gS = \Set{ x_1,x_2,x_3 } \subset S$ be a general triple of points.
The linear system $|H|$ defines an embedding $S \hto \P^7$;
let $S_0 \subset \P^4$ be the proper transform
of $S$ under the projection $\P^7 \dto \P^4$
from the plane generated by $\gS$.
We have the following:
\begin{enumerate}
\item $S_0$ has three transverse double points,
and the normalization of $S_0$ is $\Bl_\gS S$.
\item The linear system of quartics containing $S_0$ cuts out
$S_0$ as a scheme, and
defines a birational map
$\psi: \P^4 \dto \P^4$.
\item 
The inverse
$\psi^{-1}: \P^4 \dto \P^4$
is constructed the same way,
starting from another K3 surface $S'$ of degree $12$
and three points $\gS' = \Set{ x_1', x_2', x_3' }$ on $S'$.
\item We have $\Bl_{S_0}\P^4 \simeq \Bl_{S'_0}\P^4$
and it resolves $\psi$,
where $S'_0$ the proper transform
of $S'$ under the projection $\P^7 \dto \P^4$
from the plane generated by $\gS'$.
\end{enumerate}
Furthermore, if $\Pic(S) = \Z \cdot H$, 
then $S'$ is the unique Fourier--Mukai partner of $S$ such that
$S \not\simeq S'$.
Finally, if $S \hto \P^7$ and $\gS \subset S$ are defined over a subfield $\k \subset \C$,
then both $\psi$ and the pair $(S',\gS')$ are defined over $\k$.
\end{theorem}
\begin{proof}
The main statement follows from~\cite[\S1.1, Theorem 2.1]{HassettLai}.
The statement about $S \not\simeq S'$ follows from~\cite[Theorem 3.1]{HassettLai},
together with the derived invariance of Picard number for K3 surfaces 
and~\cite[Proposition 1.10]{OguisoFM}.
Finally, if $S \hto \P^7$ and $\gS \subset S$ are defined over $\k$,
then $\psi$ is defined over $\k$ by construction.
As $S_0'$ is the fundamental locus of $\psi^{-1}$,
$S_0'$ is defined over $\k$ as well.
Since $S'$ is the minimal model of the normalization of $S'_0$,
and $\gS' \subset S'$ is the image of the $(-1)$-curves,
we conclude that $(S',\gS')$ is defined over $\k$.
\end{proof}

Let $\k \subset \C$.
Let $S \subset \P^7$ be a K3 surface
of degree $12$ defined over $\k$ 
and $\Sigma \subset S$
a smooth subscheme of length three.
We say that $(S, \Sigma)$ is HL-admissible
if 
Theorem \ref{thm-HL}(1-4) holds
    for $S_\C$ and $\Sigma_\C$.

It is a construction going back to Mukai
that 
general 
K3 surfaces of degree $12$ can be obtained
as linear sections 
of a Grassmannian
$\OG_+(5,10) \subset \P^{15}$, see \cite[2.1]{HassettLai}.
Consider a general $3$-dimensional
linear section $V \subset \P^8$ over $\Q$
of this Grassmannian.
As $\OG_+(5,10)$, being a rational homogeneous variety, 
has dense $\Q$-points and since the linear section cutting off $V$ is general,
we can assume that $V$ is smooth and has a $\Q$-point. 
This implies that $V$ is $\Q$-rational~\cite[Theorem 1.1]{KuPrRatFano3}. 

Let $\P^6 \subset \P^8$ be a linear subspace
and
$\Sigma = \{ x_1, x_2, x_3 \} \subset V \cap \P^6$.
We call the 
pair 
$(\P^6, \gS)$ \emph{HL-admissible}
if the following conditions are satisfied:
\begin{enumerate}
    \item[(a)] The pencil $S_t$, $t \in \P^1$, 
    defined by intersecting
    $V$ with hyperplanes containing $\P^6$ is a Lefschetz pencil.
    \item[(b)] For one (hence for general) $t$, the pair
    $(S_t, \Sigma)$ is HL-admissible.
\end{enumerate}

\begin{lemma}
\source{motivic-invariants-birational-maps}
\notready\label{lem:HL-admis-pencil}
\source{motivic-invariants-birational-maps}
For a general $V$ as above, 
there exists an HL-admissible pair $(\P^6, \Sigma)$
defined over $\Q$.
\end{lemma}
\begin{proof} 
As both conditions (a) and (b) are nonempty Zariski open conditions,
HL-admissible pairs form a Zariski dense 
open subset $U$
in the incidence
variety 
$$\cX \cnec \Set{(x_1,x_2,x_3; L) \in V^3 \times \Gr(\P^6,\P^8)
| x_1,x_2,x_3 \in V \cap  L}.$$
It is easy to see that $\cX$ is a $\Q$-rational variety,
hence the open subset $U \subset \cX$ has (dense)
$\Q$-rational points.
\end{proof}

\begin{proposition}
\source{motivic-invariants-birational-maps}
\notready\label{pro-K3HLPic1}
\source{motivic-invariants-birational-maps}
There exist infinitely many polarized K3 surfaces $(S,H)$ of degree $12$ over $\Q$
together with three $\Q$-points $\gS = \Set{x_1,x_2,x_3} \subset S(\Q)$
such that
\begin{enumerate}
    \item $\Pic(S_\C) \simeq \Z $;
\item $(S; \gS)$ is HL-admissible. 
\end{enumerate}
The same conclusion holds if $\Q$ is replaced by any subfield $\k$ of $\C$,
and the isomorphism classes of such K3 surfaces $S$
form a set $I$ which has the cardinality of $\k$.
\end{proposition}

\begin{proof} 
The Lefschetz pencil $\{S_t\}$ 
given by the HL-admissible pair from 
Lemma~\ref{lem:HL-admis-pencil}
is defined over $\Q$,
hence Terasoma's argument in~\cite[Step 2 and Step 3 in \S3]{Terasoma}
proving~\cite[Theorem 1]{Terasoma}
applies verbatim 
and shows that there exist infinitely many 
K3 surfaces among $\{S_t\}$ defined over $\Q$
(and thus over any subfield $\k \subset \C$) with 
geometric Picard rank $1$. 
If $\k$ is uncountable,
then since ``having geometric Picard rank $1$" is a 
very general property,
the set of K3 surfaces over $\k$ 
among $\{S_t\}$ for which (1) holds
is also uncountable.


We claim that
the Lefschetz pencil
$\{S_t\}$ is non-isotrivial.
This is because the monodromy action on $H^2(S_{t,\C}, \Q)$
is nontrivial by the invariant cycle theorem,
and $\Aut(S_{t,\C}) = \{\Id\}$ if $\rho(S_{t,\C}) = 1$~\cite[Corollary 15.2.12]{HuybK3book}, 
hence
monodromy does not act by holomorphic transformations.
Thus $\{S_t\}$ contains infinitely many
isomorphism classes of K3 surfaces for any $\k \subset \C$,
uncountably many if $\k$ is uncountable.
Finally, since 
the subset of K3 surfaces $S_t$ 
in $\{S_t\}$
for which $(S_t; \gS)$ is HL-admissible  
is open, 
Proposition~\ref{pro-K3HLPic1} follows.
\end{proof}


\begin{proof}[Proof of Theorem \ref{thm:K3}]
Let $S$ be a K3 surface over $\Q$ as in Proposition~\ref{pro-K3HLPic1}
and let $S'$ be as in Theorem \ref{thm-HL}.
In particular, $S_\C \not\simeq S'_\C$ 
and $S'_\C$
is a Fourier-Mukai
partner of $S_\C$.

Let $\psi \in \Bir(\P^4)$ be the map as in Theorem~\ref{thm-HL}. 
By Theorem~\ref{thm-HL},
both $\psi$ and the K3 surface $S'$ are defined over $\Q$.
Since $S_0$ has at worst transversal double points as singularities,
the exceptional divisor of $\Bl_{S_0} : T \to \P^4$
is birational to $S_0 \times \P^1$,
which is further birational to $S \times \P^1$
by Theorem~\ref{thm-HL}(1).
Similarly, 
the exceptional divisor of $\Bl_{S'_0} : T \to \P^4$
is birational to $S' \times \P^1$.
Thus,
$$c(\psi_\k) = [S_\k \times \P^1_\k] - [S'_\k \times \P^1_\k] \in \Z[\Bir/\k].$$
Since $S_\k$ and $S_\k'$ are non-isomorphic
K3 surfaces 
because $S_\C \not\simeq S'_\C$, 
they are not stably birational, 
see e.g. Corollary~\ref{cor-CYnSB}.
Hence $c(\psi_\k) \ne 0$,
so $\psi_\k$ is motivically non-trivial.
\end{proof}



\subsection{K-trivial
threefolds}\label{ssec-IMOU}
\source{motivic-invariants-birational-maps}
\hfill

\begin{theorem}
\source{motivic-invariants-birational-maps}
\notready\label{thm:CY3}
\source{motivic-invariants-birational-maps}
Let $\k$ be an infinite field.
There exist motivically non-trivial smooth projective L-links:
\begin{equation}\label{eq:L-link-CY3}
\source{motivic-invariants-birational-maps}
 \xymatrix{
    & & T \ar[dl]_{\Bl_{Z}} \ar[dr]^{\Bl_{Z'}} & & \\
    Z \ar@{^{(}->}[r] & \XX \ar@{-->}[rr]^\psi &  & \XX' & \ar@{_{(}->}[l] Z'\\
}
\end{equation}
where $\XX$ and $\XX'$ are five-dimensional
$G_2$-Grassmannians, and
$Z$ and $Z'$ are 
K-trivial
threefolds
of Picard rank $1$.
\end{theorem}

Here, a K-trivial variety is a smooth projective variety $X$
with $\go_X \simeq \cO_X$.

\medskip

The construction of Theorem \ref{thm:CY3}
is given by Ito--Miura--Okawa--Ueda~\cite{IMOUG2} 
over a field of characteristic zero. 
We explain that the construction extends
to an arbitrary infinite field.
The proof of Theorem \ref{thm:CY3} occupies
the rest of the section.

Let $G$ be the connected and simply connected 
split simple group scheme of type $G_2$ over $\Z$.
Let $B \subset G$
be the Borel subgroup, and
$P,P' \subsetneq G$ 
the two maximal parabolic subgroups of $G$ containing $B$.
The quotients $G/B$, $G/P$, and $G/P'$ 
are smooth and projective over $\Z$ 
by~\cite[Theorem 13.33]{MilneAlgGp} 
and~\cite[Corollary XXII.5.8.5]{SGA3}.
The natural projection $p : G/B \to G/P$
is also smooth. 
Since $p_\Q : G_\Q/B_\Q \to G_\Q/P_\Q$
is a $\P^1$-fibration, so is $p : G/B \to G/P$.
Likewise, the other projection
$p' : G/B \to G/P'$ is also a smooth $\P^1$-fibration.

Our assumptions on $G$ imply that
$X^*(G) = 0$ and 
$\Pic(G) = 0$~\cite[Remark VII.1.7.a)]{RaynaudAmple},
where $X^*(G)$ is the character group of $G$.
It follows that
$$\Pic(G/P) \simeq \X^*(P) \simeq \Z$$
by~\cite[Proposition VII.1.5]{RaynaudAmple}.
Similarly, $\Pic(G/P') \simeq \Z$.
Let $M$ and $M'$ 
denote the ample generators of $\Pic(G/P)$ and $\Pic(G/P')$
respectively.
The tensor product 
$$L \cnec p^*M \otimes p'^*M'$$
is then ample as well.
Since $\deg({L_\k}_{|F_\k}) = 1$ when $\chr(\k) = 0$
where $F_\k \simeq \P^1_\k$ is a fiber of $p_\k : G_\k/B_\k \to G_\k/P_\k$
(see e.g.~\cite{IMOUG2}),
the same holds true for any field $\k$.
It follows that 
$E \cnec p_*L$ is locally free of rank $2$ and
$G/B \simeq \P(E)$ over $G/P$.

Now let $\k$ be an infinite field
and let $\gs \in H^0(G_\k/B_\k,L_\k)$ be a general section. 
Let $X_\k = Z(\gs) \subset G_\k/B_\k$ and let
$Z_\k \cnec Z(p_*\gs) \subset G_\k/P_\k$ 
with $p_*\gs$
regarded as a section of $E_\k = (p_*L)_\k$.
Similarly, let $Z'_\k \cnec Z(p'_*\gs) \subset G_\k/P'_\k$.
Since $L_\k$ is very ample~\cite[Theorem 3]{RamRam} 
and $\k$ is infinite,
$X_\k$ is smooth  
by Bertini's theorem~\cite[Th\'eor\`eme 6.3]{JouanolouBertini},~\cite[Theorem 1.10]{MukaiK31820}.
Consider the commutative diagram
\begin{equation}\label{eq:G-P-diagram}
\source{motivic-invariants-birational-maps}
\xymatrix{
    & & X_\k \ar[ddl]_{\tau} \ar[ddr]^{\tau'}
    \ar@{->}[d] & & \\
    &  & G_\k/B_\k \ar[dl]\ar[dr] &  & \\
    & G_\k/P_\k \ar@{-->}[rr]^{\psi} 
    & & G_\k/P'_\k & 
}
\end{equation}
which defines the maps $\tau$ and $\tau'$.
By construction, 
$\tau$ maps $X_\k \bss \tau^{-1}(Z_\k)$
isomorphically onto its image
and $\tau^{-1}(Z_\k) \to Z_\k$ is a $\P^1$-fibration.
By~\cite[Theorem 2]{DanilovDecomp},
$\tau$ is thus a sequence of blow ups along smooth centers
of codimension $2$ lying above $Z_\k$.
Since $X_{\k}$ has Picard rank
$\rho(X_\k) = \rho(G_\k/B_\k) = 2$
by Grothendieck--Lefschetz' theorem,
necessarily $\tau$ is the blow up of $G_\k/P_\k$ along $Z_\k$
and $Z_\k$ is connected.
Similarly, $\tau'$ is the blow up of $G_\k/P'_\k$ along $Z'_\k$,
and $Z'_\k$ is connected.
We thus have a birational map 
$$\psi \cnec \tau' \circ \tau^{-1} : G_\k/P_\k \dto G_\k/P'_\k.$$

By construction, $\psi$ is regular away from $Z_\k$
and $\psi^{-1}$ is regular away from $Z'_\k$.

\begin{lem}
\source{motivic-invariants-birational-maps}
\notready\label{lem-nSBG2}
\source{motivic-invariants-birational-maps}
$Z_\k$, $Z'_\k$ are 
K-trivial threefolds,
which are not stably birational to each other.
\end{lem}

\begin{proof}
When $\chr(\k) = 0$, Lemma~\ref{lem-nSBG2} is proven in~\cite{IMOUG2}.
Assume that $\chr(\k) = p > 0$.
Let $R$ be an integral local ring of characteristic zero with residue field $\k$
(see~\cite[Satz 20]{HasseSchmidt} or~\cite[Theorem 2.10]{HSbis} for the existence).
Since $H^j(G_\k/B_\k,L_\k)=0$ for all $j > 0$~\cite[Theorem 2]{RamRam} (and for any field $\k$),
we have
$$H^0(G_\k/B_\k,L_\k) \simeq H^0(G_R/B_R,L_R) \otimes_R \k$$
by Grauert's base change.
Thus the section $\gs \in H^0(G_\k/B_\k,L_\k)$
defining $Z_\k$
can be lifted to a section $\gs_K \in H^0(G_K/B_K,L_K)$
where $K \cnec \Frac (R)$, which is of characteristic zero.
This defines a lifting of $Z_\k$ to a smooth projective subvariety 
$Z_K \subset G_K/P_K$.  
Since $\go_{Z_K}$ is trivial, so is $\go_{Z_\k}$.
The same argument shows that $\go_{Z'_\k}$
is trivial.

To show that $Z_\k$, $Z'_\k$ are not stably birational,
we may assume $\k = \ol{\k}$. First
we show that $Z_{\k}$ and $Z'_{\k}$
are not isomorphic.
We have seen that $\rk\,\NS(X_{\k}) = 2$.
As $X_{\k} \to Z_{\k}$ and $X_{\k} \to Z'_{\k}$ are $\bP^1$-bundles,
necessarily $\rk\,\NS(Z_{\k})= \rk\,\NS(Z'_{\k}) = 1$.
Let $H$ (resp. $H'$) 
be the ample generator of $\NS(Z_{\k})$ (resp. $\NS(Z'_{\k})$).
Let $a, a' \in \Z_{>0}$ such that
$$c_1({M_{\k}}_{|Z_{\k}}) = aH \ \ \ \text{ and } \ \ \
c_1({M'_{\k}}_{|Z'_{\k}}) = a'H'.$$
As 
$$c_1({M_{\k}}_{|Z_{\k}})^3 = c_1(M_{\k})^3 \cdot Z_{\k} 
= c_1(M_K)^3 \cdot Z_K
\ \ \
\text{ and }  \ \ \ c_1({M'_{\k}}_{|Z'_{\k}})^3 = c_1(M'_{\k})^3 \cdot Z'_{\k} = c_1(M'_K)^3 \cdot Z'_K$$
are distinct and equal to 
$14$ or $42$ by~\cite{IMOUG2},
necessarily $a = a' =1$.
Hence $H^3 \ne H'^3$,
so $Z_{\k}$ is not isomorphic to $Z'_{\k}$.
It follows from Corollary~\ref{cor-CYnSB}
that $Z_\k$ is not stably birational to $Z'_\k$.
\end{proof}





\begin{proof}[Proof of Theorem \ref{thm:CY3}]
We set $\XX = G_\k/P_\k$
and $\XX' = G_\k/P'_\k$.
Both $\XX$ and $\XX'$ have a 
Bruhat decomposition
with a single cell in every dimension so that
$[\XX] = [\P^5] = [\XX']$.
Furthermore, $\XX$ and $\XX'$ are strongly
birational, because they are both
strongly rational by Lemma~\ref{lem:GP-strongly-rational}.
Thus diagram \eqref{eq:G-P-diagram} gives us the L-link
\eqref{eq:L-link-CY3}.
Finally, since  
$Z_\k$ and $Z'_\k$ are not (stably) birational by Lemma~\ref{lem-nSBG2},
the L-link~\eqref{eq:L-link-CY3} is motivically non-trivial.
\end{proof}


\begin{remark}
\source{motivic-invariants-birational-maps}
\notready
In the proof of Theorem~\ref{thm:CY3}, 
the only place where we need the assumption that
$\k$ is infinite is to guarantee 
the existence of $\gs \in H^0(G_\k/B_\k,L_\k)$
such that the vanishing loci $Z_\k \subset G_\k/P_\k$ and $Z'_\k \subset G_\k/P'_\k$
of $p_*\gs$ and $p'_*\gs$ 
are smooth. 
Based on generic smoothness,
the same conclusion of Theorem~\ref{thm:CY3} thus holds
for all but finitely many finite fields $\k$.
\end{remark}
